\documentclass[11pt]{article}
\usepackage[T1]{fontenc}
\usepackage[utf8]{inputenc}
\usepackage{lmodern,amsmath,amssymb,amsthm,microtype}
\usepackage[margin=25mm]{geometry}
\usepackage{hyperref}
\hypersetup{hidelinks,pdftitle={Strict locality and comparison over arbitrary complete fields}}
\setlength{\emergencystretch}{3em}
\newtheorem{theorem}{Theorem}
\newtheorem{lemma}[theorem]{Lemma}
\newtheorem{proposition}[theorem]{Proposition}
\newtheorem{corollary}[theorem]{Corollary}
\newcommand{\MA}{\operatorname{MA}}
\newcommand{\E}{\mathcal E}
\newcommand{\ddc}{dd^c}
\title{Strict locality and comparison\\over arbitrary complete fields}
\author{P16.4 research continuation 218\\Unreviewed partial result}
\date{7 October 2026}
\begin{document}
\maketitle
\begin{abstract}
We derive strict Borel winner locality, ordinary open-set locality
and mixed comparison for actual finite-energy non-Archimedean
potentials. The domain is a smooth geometrically integral projective
variety of any positive dimension over any complete characteristic-zero
real-valued non-Archimedean field, with an actual ample rational
bundle. The prerequisite is the companion direct finite-energy
mixed-operator and joint weighted-net theorem. Continuous locality
comes from the local uniformly approachable CLD construction.
A bounded Lipschitz test argument transports it to singular
endpoints without assuming convergence of arbitrary Borel tests
under weak convergence. Raw mass then gives comparison. The
companion prerequisite and the new proofs require independent
verification before this partial result can be accepted. No
prescribed-target equation, compactness or envelope theorem is
asserted.
\end{abstract}

\section{Actual potentials and the precise prerequisite}

Let $K$ be any complete characteristic-zero real-valued
non-Archimedean field, including trivial or nondiscrete valuation.
Let $X/K$ be smooth, geometrically integral and projective, with
$n=\dim X\geq1$, and $L$ an actual ample rational line bundle.
Fix a continuous psh reference $\psi$ and put
$T=X^{\mathrm{an}}$, $V=(L^n)>0$.
Every scalar psh potential $u$ represents the actual absolute
metric $\psi+u$, defined at all points as a decreasing-FS-net
endpoint in \cite[Definition 7.2]{BE}. Define its extended relative
energy by the infimum of continuous psh majorant energies, and
$\E^1$ by finiteness of that energy. No synthetic or abstract
potential is substituted for an endpoint.

We use exactly the following operator conclusions of the full
companion manuscript \cite{O}; until that complete prerequisite is
independently checked, the present new derivation is conditional
on it. It states that finite positive convex combinations and finite
maxima preserve the actual $\E^1$ cone, and that:

\begin{enumerate}
\item Actual $u_1,\ldots,u_n\in\E^1$ have a canonical positive
Radon mixed measure $\MA(u_1,\ldots,u_n)$ of raw mass $V$, agreeing
with the genuine continuous operator. Constants in a slot do not
change it. Continuous decreasing approximations and arbitrary
decreasing finite-energy slot nets with finite-energy endpoints give
its weak limit.
\item Every actual $w\in\E^1$ is absolutely integrable against
each such mixed output, which therefore gives zero mass to $\{w=-\infty\}$.
\item If every slot decreases to its actual $\E^1$ endpoint and
also $w_i\downarrow w$ in $\E^1$ with $w\in\E^1$, then on any
common directed index set
\begin{equation}\label{eq:prerequisite}
 \int_T w_i\,\MA(u_{1,i},\ldots,u_{n,i})
       \longrightarrow\int_T w\,\MA(u_1,\ldots,u_n).
\end{equation}
Constant weight nets are included. No locality or comparison is part
of this prerequisite.
\end{enumerate}

The companion establishes its operator from ordered continuous
Dirichlet estimates, dense scalar tests and positive polarization,
before selecting any target. Its weighted theorem uses a uniform
unbounded-test energy bound, rather than a measurable-net
monotone-convergence assertion. That full body, including every
generated dependency, is part of the verification chain needed here.
We do not replace it by a claim that a finite-energy definition alone
constructs the operator.

Finite maxima lie in $\E^1$ for an elementary reason: the actual
psh maximum is given by \cite[Corollary 7.5]{BE}, is not identically
$-\infty$, and has energy at least that of either finite-energy
operand. Constants preserve the cone by the continuous constant
law followed by its majorant infimum. These facts justify all
maxima and shifted potentials below.

\begin{theorem}\label{thm:main}
With the complete prerequisite above established, let $u,v\in\E^1$,
$m=\max(u,v)$, and $r_1,\ldots,r_{n-1}\in\E^1$.
Write $R=(r_1,\ldots,r_{n-1})$, with $R$ empty when $n=1$.
The genuine mixed operator has strict Borel winner locality:
\begin{equation}\label{eq:strict}
 \begin{split}
 \mathbf1_{\{u>v\}}\MA(m,R)&=
                  \mathbf1_{\{u>v\}}\MA(u,R),\\
 \mathbf1_{\{v>u\}}\MA(m,R)&=
                  \mathbf1_{\{v>u\}}\MA(v,R).
 \end{split}
\end{equation}
If $u=v$ at every point of an open subset $U\subset T$, then
\begin{equation}\label{eq:open-local}
                  \MA(u,R)|_U=\MA(v,R)|_U.
\end{equation}
The mixed comparison principle is
\begin{equation}\label{eq:comparison}
 \int_{\{u<v\}}\MA(v,R)\leq
                    \int_{\{u<v\}}\MA(u,R).
\end{equation}
The pure finite-energy comparison principle also holds:
\[
                  \MA(v)\{u<v\}\leq\MA(u)\{u<v\}.
\]
Its separate slotwise derivation is given below.
\end{theorem}

All winner sets are Borel: the scalar potentials are upper
semicontinuous. Differences will only be evaluated where the finitely
many involved potentials are finite. Each relevant output gives
zero mass to the excluded pole sets by the prerequisite. The literal
strict sets in \eqref{eq:strict} and \eqref{eq:comparison} include
points where only one operand is infinite, but those have zero
mass for every measure used. No arithmetic at common infinities
is required, and no countable or uncountable common full-measure
intersection across the entire net is claimed.

\section{Locality of the genuine continuous operator}

\begin{lemma}\label{lem:continuous}
For continuous fixed-$L$ psh $u,v$ and continuous remaining slots
$R$, the identities \eqref{eq:strict} hold on the ordinary open
winner sets.
\end{lemma}
\begin{proof}
The continuous psh metrics, including their maximum, are globally
uniform limits of smooth semipositive metrics by
\cite[Theorem 7.14]{BE}. On a local line-bundle trivialization their
weights are therefore locally psh-approachable in precisely
\cite[Definition 5.6.3 and Section 6.3.1]{CLD}, without needing an
assertion that all continuous CLD-psh functions are approachable.

The positive local product in \cite[Corollaries 5.6.5--5.6.6]{CLD}
is defined by uniform smooth approximation and is a morphism of
sheaves. Its complete construction checks that products on two
open sets agree on their intersection. Curvature weights transform
by the logarithm of an invertible analytic section; its Hessian
vanishes, as in \cite[Section 6.4.1]{CLD}. Hence this is the same
local metric curvature product in different trivializations.

It also agrees with the genuine global continuous measure used
here. A global uniform smooth regularization converges to the BE
measure weakly on continuous tests, by \cite[Theorem 8.4]{BE}.
On each trivializing open set the same sequence converges to the
CLD local current on compactly supported smooth tests by
Corollary 5.6.5. Every such smooth function, extended by zero,
is a global continuous test. Thus the global restriction and the
local top current have the same smooth-test values. The measure
injectivity in \cite[Proposition 5.4.6]{CLD}, proved there by local
smooth density with compact support, makes them the same Radon
measure. Analytification is boundaryless, so its entire carrier is
in the interior domain of this local calculus.

On the open set $\{u>v\}$, the local absolute weight of the maximum
equals that of $u$. The sheaf property therefore identifies the
mixed products there. Interchanging $u,v$ proves the second identity.
\end{proof}

\section{Nonlinear tests along moving finite-energy measures}

The following lemma is the passage from weak convergence and the
weighted prerequisite to the Borel cutoffs needed for locality.

\begin{lemma}\label{lem:lipschitz}
Let $\rho_i$ be finite-energy mixed outputs whose slots decrease
to actual finite-energy endpoints, and let $\rho$ be their weak
mixed limit. For $a_1,\ldots,a_q\in\E^1$, every bounded Lipschitz
$F:\mathbb R^q\to\mathbb R$ and $f\in C(T)$ satisfy
\begin{equation}\label{eq:fixed-nonlinear}
 \int_T fF(a_1,\ldots,a_q)\,d\rho_i
       \longrightarrow\int_T fF(a_1,\ldots,a_q)\,d\rho.
\end{equation}
If also $a_{s,i}\downarrow a_s$ in $\E^1$ on the same index set,
then \eqref{eq:fixed-nonlinear} holds with $a_{s,i}$ on the left.
Define $F$ to be zero outside the finite-value domain of its
arguments; this changes none of the integrals.
\end{lemma}
\begin{proof}
Choose continuous psh $a_{s,k}\downarrow a_s$ for each $s$.
For the single Radon measure $\rho$,
\[
               \int_T(a_{s,k}-a_s)\,d\rho\longrightarrow0.
\]
Here is the exact net justification. On a tail all continuous
$a_{s,k}$ have a common bound $C$, and $C-a_{s,k}$ are increasing
nonnegative continuous functions. Their strict superlevel sets
increase through a directed open cover of each superlevel of
$C-a_s$. Inner compact approximation and finite subcovers give
the equality of superlevel masses with their suprema. A common
index approximates any finite lower layer sum; truncating height
and refining layers proves convergence of the integrals, including
infinite integrals. Integrability of $a_s$ supplied by the
prerequisite makes the present integrals finite. Subtraction then
gives the displayed error limit. No arbitrary measurable-net
monotone convergence is invoked.

For each fixed continuous stage $k$, the constant-weight case of
\eqref{eq:prerequisite} and ordinary weak convergence give
\[
 \int_T(a_{s,k}-a_s)\,d\rho_i
          \longrightarrow\int_T(a_{s,k}-a_s)\,d\rho.
\]
Choose one fixed stage for each of the finitely many $s$ so that
these limiting errors are as small as desired. They are then
uniformly small on a common $i$ tail. If $M_F$ is a Lipschitz
constant for the $\ell^1$ norm, the integral difference between
the actual nonlinear test and this fixed continuous-stage test
is bounded by
\[
       \|f\|_\infty M_F
            \sum_{s=1}^q\int_T(a_{s,k}-a_s)\,d\rho_i.
\]
The analogous bound holds for $\rho$. The fixed-stage test is
continuous, so its integrals pass through weak convergence.
The two error bounds prove \eqref{eq:fixed-nonlinear}.

For moving arguments the joint weighted law and its constant-weight
case give
\[
 \int_T(a_{s,i}-a_s)\,d\rho_i
     =\int_Ta_{s,i}\,d\rho_i-\int_Ta_s\,d\rho_i
             \longrightarrow0.
\]
Every difference is nonnegative and finite on its per-index
full-measure finite-value set. The Lipschitz bound again makes
the error in replacing all moving arguments by their fixed
endpoints tend to zero. Apply \eqref{eq:fixed-nonlinear}.
\end{proof}

This lemma uses the operator's weighted theorem; weak convergence
alone would not justify the asserted Borel-test convergence.

\section{Strict Borel winner locality}

Take independent continuous psh decreasing approximations of
$u,v,r_1,\ldots,r_{n-1}$, and use their common product index $i$.
Write the approximants as $u_i,v_i,R_i$ and put
$m_i=\max(u_i,v_i)\downarrow m=\max(u,v)$.
All endpoints lie in $\E^1$. Define
\[
 \rho_i=\MA(m_i,R_i),\quad \sigma_i=\MA(u_i,R_i),\qquad
 \rho=\MA(m,R),\quad \sigma=\MA(u,R).
\]
For $\epsilon>0$ let $\chi_\epsilon(d)$ equal zero for $d\leq0$,
$d/\epsilon$ for $0<d<\epsilon$, and one for $d\geq\epsilon$.
The function $F_\epsilon(a,b)=\chi_\epsilon(a-b)$ is bounded and
Lipschitz on $\mathbb R^2$. Continuous winner locality gives
\[
 \int_T f\chi_\epsilon(u_i-v_i)\,d\rho_i
             =\int_T f\chi_\epsilon(u_i-v_i)\,d\sigma_i
                         \qquad(f\in C(T)).
\]
Apply Lemma~\ref{lem:lipschitz} separately to these two mixed-output
nets, with moving arguments $u_i,v_i$. It passes this equality to
\[
 \int_T f\chi_\epsilon(u-v)\,d\rho
             =\int_T f\chi_\epsilon(u-v)\,d\sigma.
\]
Outside the finite-value domain define the cutoff to be zero.
The two measures of all excluded pole sets are zero.

Multiplying a finite Radon measure by a bounded nonnegative Borel
function gives another finite Radon measure on this compact
Hausdorff space. Indeed compact inner and open outer approximations
to any Borel set for the original measure have weighted error
bounded by the multiplier's supremum times the original error.
Thus equality on all continuous tests identifies the two weighted
Radon measures. Let $\epsilon$ decrease along the ordinary sequence
$1/j$. On the finite-value domain the cutoffs increase to
$\mathbf1_{\{u>v\}}$. Ordinary single-measure convergence for each
of the two finite measures gives the first identity in
\eqref{eq:strict}. On the excluded pole sets both sides are zero,
so it holds for the literal winner set. Interchanging the operands
gives the second identity. There is no need to assume continuity of
a singular winner set or convergence of its indicators under weak
measure convergence.

\section{Open locality and mixed comparison}

Suppose $u=v$ on an open set $U$. For $\epsilon>0$ put
$b_\epsilon=\max(u+\epsilon,v)$.
On $U$ outside the pole sets, $u+\epsilon>v$, so strict winner
locality and constant invariance give
\[
                    \MA(b_\epsilon,R)|_U=\MA(u,R)|_U.
\]
As $\epsilon=1/j\downarrow0$, the actual finite-energy potentials
$b_\epsilon$ decrease to $m=\max(u,v)$; their mixed measures converge
weakly to $\MA(m,R)$. Test with any global continuous function
whose compact support is contained in $U$ and pass this equality
to the limit. Radon measure determination by these tests gives
$\MA(m,R)|_U=\MA(u,R)|_U$. The same argument using
$\max(u,v+\epsilon)\downarrow m$ gives equality with
$\MA(v,R)|_U$ and proves \eqref{eq:open-local}.
This claim concerns equality on an ordinary open set; equality of
potentials on an arbitrary Borel coincidence set is not asserted.

For comparison put $m_\epsilon=\max(u,v-\epsilon)$ and denote its
mixed measure by $\tau_\epsilon$. All three mixed measures have
the same raw mass $V$. Constants leave the curvature of $v$
unchanged. Strict locality consequently gives
\begin{align*}
 \MA(v,R)\{u<v-\epsilon\}
   &=\tau_\epsilon\{u<v-\epsilon\}\\
   &\leq V-\tau_\epsilon\{u>v-\epsilon\}\\
   &=V-\MA(u,R)\{u>v-\epsilon\}\\
   &=\MA(u,R)\{u\leq v-\epsilon\}\\
   &\leq\MA(u,R)\{u<v\}.
\end{align*}
All set comparisons hold on the common finite-value full-measure
set of the finitely many tests $u,v$ for each measure. Their pole
sets have zero mass, so the formulas are valid for the literal
Borel sets as written. Let $\epsilon=1/j\downarrow0$. The sets on
the first line increase to $\{u<v\}$ on the finite-value domain,
and ordinary continuity from below of the single measure
$\MA(v,R)$ proves \eqref{eq:comparison}.

The corresponding pure comparison requires a little more than
repeating fixed slots in that inequality: those slots must also
change from $u$ to $v$. Apply strict winner locality in each slot
of $\MA(m_\epsilon,\ldots,m_\epsilon)$ in turn. On
$\{u<v-\epsilon\}$ it equals $\MA(v)$; on
$\{u>v-\epsilon\}$ it equals $\MA(u)$. The same raw-mass
calculation therefore gives
\[
             \MA(v)\{u<v\}\leq\MA(u)\{u<v\}.
\]
Each intermediate slot replacement is justified by the already
proved one-slot locality with the other actual finite-energy slots
held fixed. No incorrect substitution of varying fixed slots into
\eqref{eq:comparison} is used.

\section{The unchanged original boundary}

The new derivation assumes the whole companion finite-energy
operator/weighted theorem, and additionally uses only the continuous
local CLD construction in its exact uniformly approachable domain.
The local proof is independent of a valuation's discreteness,
algebraic closure, carrier metrizability, toric invariance or measure
support. It retains the actual ample rational bundle, smooth
projectivity and characteristic zero. Raw mass is always $(L^n)$.
Its new statements need independent reconstruction; publication or
compilation alone does not accept them or their prerequisite.

Even with this comparison theorem and its full prerequisite verified,
the original P16.4 equation still requires normalized actual-potential
compactness, arbitrary-field envelope and orthogonality inputs,
both-sign energy differentiation in outer continuous directions,
the prescribed-target intrinsic finite-energy/finite-dual
characterization, actual attainment and uniqueness of every solution
modulo constants. This paper supplies none of those remaining
statements. In particular comparison is not called an existence
theorem or a proof of uniqueness, and the full original remains open.

\begin{thebibliography}{9}
\bibitem{O} P16.4 research continuation 217,
\emph{A finite-energy mixed operator over arbitrary complete fields},
complete companion manuscript, 7 October 2026. Unreviewed prerequisite;
its full body is required in the independent verification chain.
\bibitem{BE} S. Boucksom and D. Eriksson,
\emph{Spaces of norms, determinant of cohomology and Fekete points in
non-Archimedean geometry}, arXiv:1805.01016v3 (2021).
\url{https://arxiv.org/pdf/1805.01016v3}.
\bibitem{CLD} A. Chambert-Loir and A. Ducros,
\emph{Formes differentielles reelles et courants sur les espaces de
Berkovich}, arXiv:1204.6277v1 (2012).
\url{https://arxiv.org/pdf/1204.6277v1}.
\end{thebibliography}
\end{document}
